\documentclass{article}

\usepackage[spanish]{babel}
\usepackage{amsmath,amssymb}

\title{Certamen 2 - Estadística y Probabilidad}
\author{Benjamín Rivas Beltrán}
\date{\today}

\begin{document}

\maketitle

\section{Precisión brazo robótico}
Tiempos:
\[
26 - 28 - 26 - 26 - 27 - 28 - 28 - 25 - 22 - 25 - 27 - 25
\]

Se distribuyen normalmente.

\begin{itemize}
  \item Determine un intervalo de confianza del 99\% para la media de los tiempos de operación del brazo robótico.
  \item El ingeniero a cargo del proyecto sabe que si el tiempo promedio de operación del brazo robótico es mayor a 26.5 milisegundos se debe parar el proceso productivo para cambiar el brazo robótico. Con un 5\% de significancia, según la región critica, ¿qué recomienda usted?
  \item Use el valor p para confirmar su respuesta anterior.
\end{itemize}

\textbf{Solución:}
\begin{itemize}
  \item Intervalo de confianza del 99\%:
  \[
  \bar{x} = 25.9, \quad n = 12, \quad s = 1.7
  \]

  Aplicando la fórmula del intervalo de confianza:
  \[
  IC = \bar{x} \pm t_{\alpha/2, n-1} \cdot \frac{s}{\sqrt{n}}
  \]
  Donde:
  \[
  t_{\alpha/2, n-1} = t_{0.005, 11} \approx 3.106
  \]
  Entonces:
  \[
  IC = 25.9 \pm 3.106 \cdot \frac{1.7}{\sqrt{12}} \approx 25.9 \pm 1.52
  \]
  Por lo tanto, el intervalo de confianza del 99\% es aproximadamente:
  \[
  (24.38, 27.42)
  \]

  \item Prueba de hipótesis:
  Las hipótesis son:
  \[
  H_0: \mu \leq 26.5 \quad \text{} H_1: \mu > 26.5
  \]

  Calculamos el estadístico de prueba:
  \[
  t = \frac{\bar{x} - \mu_0}{s/\sqrt{n}} = \frac{25.9 - 26.5}{1.7/\sqrt{12}} \approx -1.22
  \]
  El valor crítico para un nivel de significancia del 5\% y 11 grados de libertad es:
  \[
  t_{0.05, 11} = 1.796
  \]

  Dado que \( t = -1.22 < 1.796 \), no rechazamos \( H_0 \). Por lo tanto, no hay suficiente evidencia para recomendar detener el proceso productivo.

  \item Valor p:
  El valor p se calcula como:
  \[
  p = P(T > -1.22) \approx 0.88
  \]
  Dado que \( p = 0.88 > 0.05 \), no rechazamos \( H_0 \). Esto confirma nuestra decisión anterior de no detener el proceso productivo. 
\end{itemize}

\section{Prueba ingenieros civiles}
Los datos se distribuyen normalmente con una media de 580 puntos y una desviación estándar de 38 puntos. El 6.82\% de los mejores puntajes son seleccionados, los que tienen entre 595 y 620 quedan en lista de espera.

\begin{itemize}
  \item Determine el puntaje mínimo que debe obtener un ingeniero civil para ser seleccionado.
  \item Si la prueba la rinden 40 ingenieros civiles, ¿cuántos postulantes quedarían en lista de espera?
  \item Determine e interprete el percentil 85.
\end{itemize}

\textbf{Solución:}
\begin{itemize}
  \item Puntaje mínimo para ser seleccionado:
  Sea \( X \) la variable aleatoria que representa los puntajes de los ingenieros civiles. Sabemos que \( X \sim N(580, 38^2) \).

  Queremos encontrar el puntaje mínimo \( x \) tal que:
  \[
  P(X > x) = 0.0682
  \]
  Esto implica que:
  \[
  P(X \leq x) = 1 - 0.0682 = 0.9318
  \]
  
  Estandardizamos la variable:
  \[
  Z = \frac{X - 580}{38}
  \]
  Buscamos \( z \) tal que \( P(Z \leq z) = 0.9318 \). Usando la tabla de la distribución normal estándar, encontramos que \( z \approx 1.48 \).

  Entonces:
  \[
  \frac{x - 580}{38} = 1.48 \implies x = 580 + 1.48 \cdot 38 \approx 580 + 56.24 \approx 636.24
  \]
  Por lo tanto, el puntaje mínimo para ser seleccionado es aproximadamente 636 puntos. 

  \item Número de postulantes en lista de espera:
  Queremos encontrar el número de ingenieros civiles que quedan en lista de espera, es decir, aquellos con puntajes entre 595 y 620.
  Calculamos las probabilidades:
  \begin{align*}
  P(595 \leq X \leq 620) = P\left(\frac{595 - 580}{38} \leq Z \leq \frac{620 - 580}{38}\right) \\ = P(0.395 \leq Z \le 1.053) = P(Z \leq 1.053) - P(Z \leq 0.395)
  \end{align*}

  Reemplazando los valores:
  \[
  P(Z \leq 1.053) \approx 0.853, \quad P(Z \leq 0.395) \approx 0.653
  \]
  Entonces:
  \[
  P(595 \leq X \leq 620) \approx 0.853 - 0.653 = 0.2
  \]
  Si 40 ingenieros civiles rinden la prueba, el   número esperado de postulantes en lista de espera es:
  \[
  40 \cdot 0.2 = 8
  \]

  \item Percentil 85:
  El percentil 85 es el valor \( x \) tal que:
  \[
  P(X \leq x) = 0.85
  \]
  Estandardizamos:
  \[
  Z = \frac{x - 580}{38}
  \]
  Buscamos \( z \) tal que \( P(Z \leq z) = 0.85 \). Usando la tabla de la distribución normal estándar, encontramos que \( z \approx 1.03 \).

  Entonces:
  \[
  \frac{x - 580}{38} = 1.03 \implies x = 580 + 1.03 \cdot 38 \approx 580 + 39.14 \approx 619.14
  \]
  Por lo tanto, el percentil 85 es aproximadamente 619 puntos. Esto significa que el 85\% de los ingenieros civiles obtuvieron un puntaje menor o igual a 619 puntos, y el 15\% obtuvo un puntaje mayor a este valor.
\end{itemize}

\section{Martillo demoledor de cncreta}
Vida útil de 4 años, con distribución exponencial.

\begin{itemize}
  \item Obtenga la función de densidad de probabilidad y la función de distribución acumulada.
  \item Determine la probabilidad de que un martillo demoledor dure a lo más 7 años.
  \item Si el martillo lleva funcionando más de 4 años, ¿cuál es la probabilidad de que dure a lo más 7 años?
  \item Determine e interprete el quintil 4.
\end{itemize}

\textbf{Solución:}
\begin{itemize}
  \item Función de densidad de probabilidad y función de distribución acumulada:
  Dado que la vida útil del martillo demoledor sigue una distribución exponencial con media \( \mu = 4 \) años, la tasa \( \lambda \) se calcula como:
  \[
  \lambda = \frac{1}{\mu} = \frac{1}{4} = 0.25
  \]
  La función de densidad de probabilidad es:
  \[
  f(x) = \lambda e^{-\lambda x} = 0.25 e^{-0.25 x}, \quad x \geq 0
  \]
  La función de distribución acumulada es:
  \[
  F(x) = 1 - e^{-\lambda x} = 1 - e^{-0.25 x}, \quad x \geq 0
  \]

  \item Probabilidad de que un martillo demoledor dure a lo más 7 años:
  \[
  P(X \leq 7) = F(7) = 1 - e^{-0.25 \cdot 7} \approx 1 - e^{-1.75} \approx 1 - 0.1738 \approx 0.8262
  \]

  La probabilidad de que un martillo demoledor dure a lo más 7 años es aproximadamente 0.8262 o 82.62\%.

  \item Probabilidad de que dure a lo más 7 años dado que lleva funcionando más de 4 años:
  Usamos la fórmula de probabilidad condicional:
  \[
  P(X \leq 7 \mid X > 4) = \frac{P(4 < X \leq 7)}{P(X > 4)} = \frac{F(7) - F(4)}{1 - F(4)}
  \]
  Calculamos \( F(4) \):
  \[
  F(4) = 1 - e^{-0.25 \cdot 4} = 1 - e^{-1} \approx 1 - 0.3679 \approx 0.6321
  \] 
  Calculamos \( F(7) \):
  \[
  F(7) \approx 0.8262
  \]
  Entonces:
  \[
  P(X \leq 7 \mid X > 4) = \frac{0.8262 - 0.6321}{1 - 0.6321} = \frac{0.1941}{0.3679} \approx 0.5277
  \]
  La probabilidad de que un martillo demoledor dure a lo más 7 años dado que lleva funcionando más de 4 años es aproximadamente 0.5277 o 52.77\%.

  \item Quintil 4:
  El quintil 4 es el valor \( x \) tal que:
  \[
  P(X \leq x) = 0.8
  \]
  Usamos la función de distribución acumulada:
  \[
  0.8 = 1 - e^{-0.25 x} \implies e^{-0.25 x} = 0.2 \implies -0.25 x = \ln(0.2) \implies x = -\frac{\ln(0.2)}{0.25} \approx 6.4378
  \]
  Por lo tanto, el quintil 4 es aproximadamente 6.44 años. Esto significa que el 80\% de los martillos demoledores tienen una vida útil menor o igual a 6.44 años, y el 20\% restante tiene una vida útil mayor a este valor.
\end{itemize}

\section{Ladrillos construcción}
Los pesos de los ladrillos de construcción se distribuyen normalmente con una media de 4.8 kg, y se extrae una muestra de 36 ladrillos con peso promedio de 4.72 kg y desviación estándar de 0.2 kg.

\begin{itemize}
  \item Con un 5\% de significancia, determine si el peso promedio de los ladrillos es menor a 4.8 kg. Utilice la región crítica y el valor p.
  \item Determine un intervalo de confianza del 95\% para el peso promedio de los ladrillos.
\end{itemize}

\begin{itemize}
  \item Prueba de hipótesis:
  Las hipótesis son:
  \[
  H_0: \mu = 4.8 \quad \text{} H_1: \mu < 4.8
  \]

  Dado que la muestra es grande (\( n = 36 \)), podemos usar la distribución normal para el estadístico de prueba:
  \[
  z = \frac{\bar{x} - \mu_0}{s/\sqrt{n}} = \frac{4.72 - 4.8}{0.2/\sqrt{36}} = \frac{-0.08}{0.0333} = -2.4
  \]

  El valor crítico para un nivel de significancia del 5\% en una prueba unilateral es:
  \[
  z_{0.05} = -1.645
  \]

  Dado que \( z = -2.4 < -1.645 \), rechazamos \( H_0 \). Por lo tanto, hay suficiente evidencia para concluir que el peso promedio de los ladrillos es menor a 4.8 kg.

  Esta decisión se confirma con el valor p:
  \[
  p = P(Z < -2.4) \approx 0.0082
  \]
  Dado que \( p = 0.0082 < 0.05 \), rechazamos \( H_0 \) y concluimos que el peso promedio de los ladrillos es menor a 4.8 kg.

  \item Intervalo de confianza del 95\%:
  El intervalo de confianza del 95\% para la media se calcula como:
  \[
  IC = \bar{x} \pm z_{\alpha/2} \cdot \frac{s}{\sqrt{n}}
  \]
  Donde \( z_{\alpha/2} = z_{0.025} \approx 1.96 \).
  \[
  IC = 4.72 \pm 1.96 \cdot \frac{0.2}{\sqrt{36}} = 4.72 \pm 1.96 \cdot 0.0333 = 4.72 \pm 0.0653
  \]
  Por lo tanto, el intervalo de confianza del 95\% para el peso promedio de los ladrillos es:
  \[
  (4.6547, 4.7853)
  \]
  
\end{itemize}

\end{document}